%%=============================================================================
%% Inleiding
%%=============================================================================

\chapter{\IfLanguageName{dutch}{Inleiding}{Introduction}}%
\label{ch:inleiding}


<<<<<<< HEAD
Het doel van dit onderzoek is om te achterhalen hoe de ontwikkeling van AI-gebaseerde cheats ontwikkeld worden en hoe we dit kunnen tegen gaan en hiervoor is de volgende onderzoeksvraag opgesteld: Hoe worden AI-gebaseerde cheats ontwikkeld en welke methoden kunnen worden ontwikkeld deze te identificeren op basis van spelersgedrag en interactiepatronen? Spelersgedrag is hierbij de fysieke muisbeweging en reactiesnelheid waar de speler op bepaalde instanties reageert.
\\\\
Om een antwoord te kunnen geven op de onderzoeksvraag zullen we eerst kijken hoe het vervalsen van muis en toetsenbord input gebeurd door een 2de computer. Vervolgens wordt er onderzocht hoe een model getrained wordt op het detecteren en reageren op anomalieën. Ten slotte wordt een analyse gemaakt op de verzamelde gebruikers input om anomalieën te detecteren die zouden kunnen lijden naar het gebruiken van AI-gebaseerde cheats. 
\\\\
De resultaten van dit onderzoek zullen leiden tot een raamwerk dat de specifieke gedrags- en inputkenmerken identificeert die wijzen op het gebruik van AI-gebaseerde cheats. Dit zal resulteren in een aanbeveling voor de meest geschikte machine learning-methoden voor het detecteren van afwijkend spelersgedrag, waarmee een oplossing voor dit probleem kan gevonden worden.
=======
Het doel van dit onderzoek is om te achterhalen hoe de ontwikkeling van AI-gebaseerde cheats ontwikkeld wordt en hoe dit kan worden tegen gaan.
Hiervoor is de volgende onderzoeksvraag opgesteld: Hoe worden AI-gebaseerde cheats ontwikkeld en welke methodes worden er gebruikt om deze te identificeren aan de hand van een analyse op spelersgedrag en interactiepatronen.
Spelersgedrag is hierbij de fysieke muisbeweging en reactiesnelheid waar de speler op bepaalde instanties reageert.
\\\\
Om een antwoord te kunnen geven op de onderzoeksvraag wordt er eerst gekeken naar hoe het vervalsen van muis- en toetsenbordinput gebeurd door een 2de computer. Vervolgens wordt er onderzocht hoe een model getrained kan worden op het detecteren en reageren op anomalieën (vijandige spelers). Ten slotte wordt een analyse gemaakt op de verzamelde gebruikers input om anomalieën te detecteren die zouden lijden naar een gebruiken van AI-gebaseerde cheats. 
\\\\
De resultaten van dit onderzoek leiden tot een raamwerk dat de specifieke gedrags- en inputkenmerken identificeert die wijzen op het gebruik van AI-gebaseerde cheats. Dit resulteert in een aanbeveling voor de meest geschikte machine learning-methoden voor het detecteren van afwijkend spelersgedrag, waarmee een oplossing voor dit probleem kan gevonden worden.
\\\\


>>>>>>> 1cb6c8cf8372074a7f79583067454376aa80bc2b

\section{\IfLanguageName{dutch}{Probleemstelling}{Problem Statement}}%
\label{sec:probleemstelling}

<<<<<<< HEAD
Met de groeiende opkomst van AI-gebaseerde cheats worden traditionele anti-cheatmethoden steeds minder effectief. Deze moderne cheats kunnen menselijke input realistisch nabootsen of manipuleren via een tweede computer, waardoor klassieke detectiesystemen moeite hebben om afwijkend gedrag te herkennen. Dit vormt een dringend probleem voor organisaties die betrokken zijn bij het waarborgen van eerlijk spel.  
=======
Door de groeiende opkomst van AI-gebaseerde cheats worden traditionele anti-cheatmethoden steeds minder effectief. Deze moderne cheats kunnen menselijk verwachte input vervalsen via een tweede computer, waardoor klassieke detectiesystemen moeite hebben om afwijkend gedrag te herkennen. Dit vormt een dringend probleem voor organisaties die betrokken zijn bij het waarborgen van eerlijk spel.  
>>>>>>> 1cb6c8cf8372074a7f79583067454376aa80bc2b

De concrete doelgroep van dit onderzoek bestaat uit game-ontwikkelaars, anti-cheatteams en onderzoekers die zich bezighouden met gedragsanalyse in online gamingomgevingen. Zij hebben nood aan nieuwe benaderingen die zich richten op het detecteren van subtiele anomalieën in input- en gedragsdata die specifiek verband houden met AI-gebaseerde cheats.

\section{\IfLanguageName{dutch}{Onderzoeksvraag}{Research Question}}%
\label{sec:onderzoeksvraag}

De centrale onderzoeksvraag van deze bachelorproef luidt als volgt:

\begin{quote}
<<<<<<< HEAD
\textbf{Hoe worden AI-gebaseerde cheats ontwikkeld en welke methoden kunnen worden ontworpen om deze te identificeren op basis van spelersgedrag en interactiepatronen?}
=======
\textbf{Hoe worden AI-gebaseerde cheats ontwikkeld en welke methodes worden er gebruikt om deze te identificeren aan de hand van een analyse op spelersgedrag en interactiepatronen.}
>>>>>>> 1cb6c8cf8372074a7f79583067454376aa80bc2b
\end{quote}

Deze vraag wordt verder opgesplitst in de volgende deelvragen:

\begin{itemize}
<<<<<<< HEAD
=======
    \item Welke modules zijn er nodig om een effectieve cheat te maken?
>>>>>>> 1cb6c8cf8372074a7f79583067454376aa80bc2b
    \item Hoe worden muis- en toetsenbordinputs gemanipuleerd met behulp van een tweede computer?
    \item Welke machine-learningmethoden zijn geschikt om anomalieën in spelersgedrag te detecteren?
    \item Welke gedrags- en inputpatronen zijn kenmerkend voor het gebruik van AI-gebaseerde cheats?
    \item Welke gegevens leveren de meest betrouwbare signalen voor een detectiesysteem?
\end{itemize}

\section{\IfLanguageName{dutch}{Onderzoeksdoelstelling}{Research Objective}}%
\label{sec:onderzoeksdoelstelling}

<<<<<<< HEAD
Het doel van deze bachelorproef is het ontwikkelen van een raamwerk dat gedrags- en inputkenmerken identificeert die wijzen op de aanwezigheid van AI-gebaseerde cheats. Dit gebeurt door het analyseren van hoe AI-cheats worden opgebouwd en hoe zij gebruikersinput manipuleren, het trainen van een model voor anomaliedetectie, en het evalueren van methoden voor het herkennen van afwijkend spelersgedrag.

Het beoogde resultaat is een proof-of-concept en een onderbouwde aanbeveling voor detectiemethoden die toepasbaar zijn binnen moderne anti-cheatsystemen.

=======
Het doel van deze bachelorproef is het ontwikkelen van een raamwerk dat gedrags- en inputkenmerken identificeert die wijzen op de aanwezigheid van AI-gebaseerde cheats. 
Dit gebeurt door het analyseren van hoe AI-cheats worden opgebouwd en hoe zij gebruikersinput manipuleren, het trainen van een model voor anomaliedetectie, en het evalueren van methoden voor het herkennen van afwijkend spelersgedrag.
\\\\
Het verwachte resultaat is een proof-of-concept en een onderbouwde aanbeveling voor detectiemethoden die toepasbaar zijn binnen moderne anti-cheatsystemen.

>>>>>>> 1cb6c8cf8372074a7f79583067454376aa80bc2b
\section{\IfLanguageName{dutch}{Opzet van deze bachelorproef}{Structure of this Bachelor Thesis}}%
\label{sec:opzet-bachelorproef}

Deze bachelorproef is als volgt opgebouwd:

\begin{itemize}
<<<<<<< HEAD
    \item \textbf{Hoofdstuk 1 – Inleiding:} bespreekt het probleem en de relevantie van cheatdetectie.
    \item \textbf{Hoofdstuk 2 – Analyse van inputmanipulatie:} onderzoekt hoe een tweede computer muis- en toetsenbordinput kan vervalsen.
    \item \textbf{Hoofdstuk 3 – Modeltraining:} bespreekt hoe anomaliedetectiemodellen worden getraind.
    \item \textbf{Hoofdstuk 4 – Data-analyse:} analyseert verzamelde inputdata om afwijkend gedrag te identificeren.
    \item \textbf{Hoofdstuk 5 – Resultaten:} presenteert het raamwerk en de aanbevolen methoden.
    \item \textbf{Hoofdstuk 6 – Conclusie en aanbevelingen.}
=======
    \item \textbf{Hoofdstuk 1 – Vaktermen binnen FPS games:} Dit hoofdstuk legt de belangrijkste terminologie uit die gebruikt wordt binnen FPS-games, zodat de context van latere technische concepten correct begrepen wordt.
    
    \item \textbf{Hoofdstuk 2 – Recoilcompensatie:} Hier wordt uitgelegd hoe recoil in FPS-games werkt en op welke manieren spelers of systemen deze kunnen compenseren, zowel handmatig als geautomatiseerd.
    
    \item \textbf{Hoofdstuk 3 – AI gebaseerde cheats:} Dit hoofdstuk bespreekt de werking van AI-gestuurde cheats en hoe deze input kunnen automatiseren om een voordeel te verkrijgen in competitieve games.
    
    \item \textbf{Hoofdstuk 4 – Trainen van het model:} In dit hoofdstuk wordt toegelicht hoe een machine learning model wordt opgebouwd, getraind en geëvalueerd op basis van verzamelde gedragsdata.
    
    \item \textbf{Hoofdstuk 5 – Fysieke proof-of-concept:} Dit hoofdstuk beschrijft de hardwarematige implementatie van het systeem, waarbij een fysieke setup wordt gebruikt om AI-input naar een realistische omgeving te vertalen.
    
    \item \textbf{Hoofdstuk 6 – Integratie van alle componenten:} Hier wordt uitgelegd hoe de verschillende software- en hardwarecomponenten samengevoegd worden tot één werkend systeem.
    
    \item \textbf{Hoofdstuk 7 – Detectie van valsspelen via machine learning:} Dit hoofdstuk behandelt de ontwikkeling van het anti-cheatsysteem dat gedragsanalyse en machine learning gebruikt om AI-gegenereerde input te detecteren.
>>>>>>> 1cb6c8cf8372074a7f79583067454376aa80bc2b
\end{itemize}

% Het is gebruikelijk aan het einde van de inleiding een overzicht te
% geven van de opbouw van de rest van de tekst. Deze sectie bevat al een aanzet
% die je kan aanvullen/aanpassen in functie van je eigen tekst.
